% discussion.tex — Paper 6, Section VII: Discussion
% ASSERT_CONVENTION: natural_units=natural, metric_signature=mostly_plus, entropy_base=nats, modular_hamiltonian=K_minus_ln_rho, coupling_convention=H_sum_hxy, state_normalization=Tr_rho_1
\section{Discussion}
\label{sec:discussion}

\subsection{What is derived versus what is input}
\label{sec:derived-vs-input}

The derivation chain summarized in Table~\ref{tab:chain} establishes
Einstein's field equations from self-modeling through nine links.  It is
important to state precisely what is derived, what is input, and what is
standard methodology.

\emph{Derived from self-modeling.}---The local algebra
$M_n(\mathbb{C})^{sa}$ with the L\"uders sequential product is forced by
faithful self-modeling~\cite{Paper5}. The SWAP interaction Hamiltonian
$H = \sum_{\langle x,y\rangle} J\,\swapop$ is forced by diagonal $U(n)$
covariance and Schur--Weyl duality (Sec.~\ref{sec:lattice}).  The
sub-volume entanglement structure follows from the locality of this
Hamiltonian via the WVCH mutual information
bound~\cite{WVCH2008}, known ground-state entanglement properties of the
Heisenberg model, and the entanglement first law combined with modular
Hamiltonian locality
(Sec.~\ref{sec:arealaw}).  The entanglement first law
$\delta\ent = \delta\langle\modH\rangle$ is an exact identity of quantum
information theory, requiring no additional assumptions.

\emph{Input: lattice topology.}---The graph $G = (V, E)$ defining which
sites are nearest neighbors is not derived from self-modeling.  The paper
derives the emergent \emph{metric}---the assignment of distances
compatible with the entanglement structure---but not the
\emph{topology}.  The spatial dimension is set by the graph structure:
a $d$-dimensional regular lattice yields $d$ emergent spatial dimensions.
This is an irreducible input to the framework.

\emph{Input: antiferromagnetic coupling.}---The sign of~$J$ is not
determined by self-modeling.  Only $J > 0$ (antiferromagnetic) supports
nontrivial entanglement; the ferromagnetic ground state is a product
state with zero entanglement (Sec.~\ref{sec:lattice}).  We work with
$J > 0$ throughout; this is an input.

\emph{Standard methodology: Wilsonian continuum limit.}---The passage
from a lattice theory to a smooth Riemannian manifold at long wavelengths
$L \gg a$ (where $a$ is the lattice spacing) employs a Wilsonian
universality argument.  This is standard methodology in lattice QCD,
causal dynamical triangulations, and Regge calculus.  However, its
applicability to the self-modeling Hamiltonian in $d \geq 2$ is not
guaranteed---see Gap~1 below.

\emph{Standard: Jacobson (1995) thermodynamic argument.}---The bridge
from area-law entropy to Einstein's equation uses the Clausius relation
$\delta Q = T\,\delta\ent$ on local Rindler horizons, following
Jacobson~\cite{Jacobson1995}.  This argument is UV-agnostic: it applies
to any UV completion that provides $\ent = \eta\,\mathcal{A}$ and
emergent Lorentz invariance.  We provide the specific UV completion;
Jacobson provides the bridge to GR.

\subsection{Remaining gaps}
\label{sec:gaps}

Four gaps remain in the derivation chain.

\paragraph{Gap 1: Emergent Lorentz invariance.}
The lattice has a preferred frame.  The Jacobson argument requires local
Lorentz invariance to construct Rindler horizons and apply the Unruh
effect.  In $d = 1$, the infrared theory is the $SU(2)_1$ WZW CFT,
which is exactly Lorentz-invariant.  In $d \geq 2$, the Heisenberg
antiferromagnet is expected to exhibit N\'eel order, with spin-wave
excitations that see an emergent Lorentz symmetry at long wavelengths.
This is standard condensed matter lore but not rigorously proved.
This gap replaces the old conformal approximation and tensoriality
gaps---it is a single, cleaner gap that asks a standard condensed matter
question rather than a bespoke one.

\paragraph{Gap 2: Lattice Bisognano--Wichmann.}
The Jacobson argument uses the Unruh temperature $T = \hbar\kappa/(2\pi)$,
which follows from the Bisognano--Wichmann theorem in continuum QFT.  On
the lattice, we need the modular flow of the ground state restricted to a
half-chain to approximate a local boost near the entangling surface.
Our numerical evidence is \emph{consistent} with this: the short-range
fraction of $\modH$ is $\mathrm{SRF} = 0.9993$ for the Heisenberg AFM
at $N = 16$ (Sec.~\ref{sec:ed-ka}), indicating that $\modH$ is
concentrated at the boundary.  However, boundary concentration is
\emph{necessary} but not \emph{sufficient} for the BW property: many
short-range operators localized at the boundary are not boost generators.
The full BW claim requires that the modular flow acts geometrically as a
local boost, not merely that it is short-ranged.  Establishing this on
the lattice remains open.  This is nonetheless a weaker requirement than
the full conformal modular Hamiltonian used in Jacobson's 2016
argument~\cite{Jacobson2016}: we need only that the modular flow
\emph{near the boundary} approximates a local boost, not that it takes
the exact Casini--Huerta--Myers form~\cite{CasiniHuertaMyers2011}
everywhere.

\paragraph{Gap 3: Local equilibrium.}
The Raychaudhuri step assumes $\theta = \sigma_{ab} = 0$ at the
bifurcation surface of the local Rindler horizon.  This is physically
natural for the self-modeling fixed point: the ground state is the
equilibrium state of the self-modeling dynamics, so its entanglement
structure is stationary.  In the continuum limit, stationarity of the
vacuum entanglement across a Rindler horizon translates to vanishing
expansion and shear at the bifurcation surface.

\paragraph{Gap 4: Continuum limit.}
The Wilsonian argument that the lattice produces a smooth Riemannian
manifold $(M, \gab)$ at long wavelengths is a physical argument, not a
rigorous construction.  A constructive proof would require showing that
lattice correlation functions converge (in an appropriate sense) to those
of a local quantum field theory on a Riemannian manifold as $a \to 0$.
This is a hard open problem shared by \emph{all} lattice quantum gravity
programs---lattice QCD, causal dynamical triangulations, Regge
calculus---and is not specific to self-modeling.  The infrared behavior
of the $d \geq 2$ Heisenberg model (likely N\'eel-ordered rather than
conformal) means the emergent geometry may carry specific
signatures---e.g., a mass gap corresponding to a cosmological
constant---rather than being exactly Minkowski.  This is potentially
interesting rather than problematic.

\emph{Note on eliminated assumptions.}---Three features of the prior
version of this paper are eliminated in the present derivation.  (i)~The
maximal vacuum entanglement hypothesis (MVEH)---previously an input
reframed as a structural identification via the Connes--Rovelli thermal
time hypothesis---is no longer needed.  Jacobson
(1995)~\cite{Jacobson1995} uses the Clausius relation on local horizons,
not an entanglement equilibrium condition.  (ii)~The conformal modular
Hamiltonian, which Jacobson (2016)~\cite{Jacobson2016} required to
evaluate $\delta\ent_{\mathrm{mat}}$, is replaced by the simpler
$\ent = \eta\,\mathcal{A}$ input.  (iii)~The tensoriality assumption of
the Lovelock route is eliminated: Jacobson's null-direction argument
extracts the tensor equation directly.

\subsection{Comparison with related work}
\label{sec:comparison}

\paragraph{Jacobson (1995)~\cite{Jacobson1995} and (2016)~\cite{Jacobson2016}.}
Jacobson's 1995 derivation treats the Einstein equation as an equation of
state: $\delta Q = T\,\delta\ent$ on local Rindler horizons, with the
Unruh temperature and Bekenstein--Hawking entropy, yields Einstein's
equation.  The inputs are the proportionality $\ent = \eta\,\mathcal{A}$,
the Unruh effect, and the Clausius relation.  His 2016 entanglement
equilibrium argument~\cite{Jacobson2016} is more refined, decomposing
the entropy into UV and matter sectors and using MVEH, but requires a
conformal modular Hamiltonian.  We use the 1995 argument because it is
simpler and requires fewer assumptions in $d \geq 2$.  Our contribution
is the UV completion upstream of Jacobson's inputs: we \emph{derive} the
area-law entropy $\ent = \eta\,\mathcal{A}$ from the self-modeling
lattice (Links~L1--L5 in Table~\ref{tab:chain}), rather than assuming
it.  The UV content (L1--L5) feeds equally well into either Jacobson
argument.

\paragraph{Cao--Carroll--Michalakis (2017)~\cite{CCM2017}.}
CCM derive the spatial constraint equation (the Hamiltonian constraint of
general relativity) from the entanglement structure of a
finite-dimensional Hilbert space.  Our result gives the full
\emph{spacetime} Einstein equation, including the dynamical sector, not
only the spatial constraint.  The spatial part of our Einstein equation is
consistent with CCM's constraint.  Both approaches share the insight that
entanglement structure encodes geometry; the key difference is that CCM
reconstruct spatial geometry from entanglement, while we (following
Jacobson) derive the full gravitational dynamics from the thermodynamics
of local Rindler horizons.

\paragraph{Lashkari--McDermott--Van Raamsdonk (2014)~\cite{LMVR2014}.}
LMVR showed that the entanglement first law in holographic conformal field
theories implies the \emph{linearized} Einstein equation.  Our
derivation, following Jacobson~\cite{Jacobson1995}, extends this to the
full \emph{nonlinear} Einstein equation via the Clausius relation.
The LMVR result is recovered as the linearized limit of our
Eq.~(\ref{eq:einstein}) around flat space.  Crucially, the LMVR derivation
assumes AdS/CFT, while ours is non-holographic.

\paragraph{Faulkner \textit{et~al.}\ (2014)~\cite{Faulkner2014}.}
Faulkner~\textit{et~al.}\ derived the full nonlinear Einstein equations
from entanglement in the holographic setting, extending LMVR beyond the
linearized regime.  Their derivation, like LMVR's, assumes the AdS/CFT
correspondence.  Our approach is non-holographic: the self-modeling
lattice does not assume a bulk/boundary duality.

\subsection{Connection to Papers~5 and~7}
\label{sec:paper5-connection}

Paper~5~\cite{Paper5} showed that faithful self-modeling forces quantum
mechanics: the state space is $M_n(\mathbb{C})^{sa}$ with the L\"uders
sequential product and the conjugate-transpose involution.  This paper
shows that self-modeling forces general relativity: the locality of
self-modeling produces area-law entanglement, and the Jacobson
thermodynamic argument converts this to Einstein's equations.

Together, the two papers establish that the self-modeling
principle---combined with a lattice topology---yields both quantum
mechanics and general relativity as necessary consequences.  Quantum
mechanics emerges from the algebraic structure of single-site
self-modeling (the $B$--$M$ boundary interaction).  General relativity
emerges from the entanglement structure of the multi-site lattice (the
inter-site boundary interactions).

A natural question is whether the Standard Model gauge group also follows
from self-modeling.  If so, the spectral action on the exceptional Jordan
algebra $h_3(\mathbb{O})$ may unify the gravitational and
gauge-theoretic derivations into a single geometric framework (future
work).

\subsection{What this paper does not claim}
\label{sec:not-claimed}

We state explicitly what lies outside the scope of this paper:
\begin{itemize}
\item \textbf{Newton's constant.}  The value $\Gnewton = 1/(4\eta)$
  depends on the entanglement entropy density~$\eta$, which in turn
  depends on the lattice local dimension~$n$ and the lattice spacing~$a$.
  These are not determined by self-modeling.
\item \textbf{Spacetime dimension.}  The spatial dimension~$d$ is set by
  the lattice topology $G = (V, E)$, which is an irreducible input.
  Self-modeling does not select $d = 3$.
\item \textbf{Cosmological constant.}  $\Lambda$ appears as an
  undetermined integration constant in the Jacobson
  derivation~\cite{Jacobson1995, JacobsonSperanza2019}.  We do not
  predict its value.
\item \textbf{Null energy condition.}  The sign chain producing
  attractive gravity (Sec.~\ref{sec:jacobson95}) assumes the NEC
  ($\Tab\, k^a k^b \geq 0$ for null $k^a$).  The tensor
  equation~\eqref{eq:einstein} holds without the NEC; only the sign
  interpretation requires it.
\item \textbf{Emergent Lorentz invariance.}  We do not prove that
  the Heisenberg lattice produces Lorentz-invariant long-wavelength
  physics in $d \geq 2$.  This is stated as Gap~1.
\end{itemize}

\subsection{Future work}
\label{sec:future}

Several directions are natural extensions of this work.

\emph{Spatial dimension from self-modeling.}---Can constraints on the
lattice topology $G = (V, E)$ be derived from self-modeling requirements?
If so, this would reduce the input set further, potentially selecting
$d = 3$ spatial dimensions.

\emph{Matter content and spectral action.}---The self-modeling lattice
produces a gravitational sector.  The connection to the Standard Model
gauge group via the exceptional Jordan algebra $h_3(\mathbb{O})$ and
spectral geometry is explored in a companion paper.  If the spectral
action on the observer's algebra reproduces both Einstein gravity and the
SM gauge group, this would unify the present derivation with the gauge
theory derivation.

\emph{Spectral gap.}---Proving that the self-modeling Hamiltonian has a
spectral gap above the ground state (for appropriate sign of~$J$ and
appropriate dimension) would strengthen the area-law argument, bringing
the Hastings theorem~\cite{Hastings2007} and the Brand\~ao--Horodecki
result~\cite{BrandaoHorodecki2013} to bear directly.

\emph{Larger numerical simulations.}---The exact diagonalization results
in Sec.~\ref{sec:numerical} are limited to $N \leq 20$ sites.  Tensor
network methods (DMRG, TEBD) could push to $N \sim 10^2$--$10^3$ in one
dimension, providing stronger evidence for area-law scaling and lattice
Bisognano--Wichmann convergence.
